\documentclass[11pt]{article}
% BEGIN SHARED COURSE STYLE
% Shared Math 615 style. Embedded verbatim in standalone published sources.
\RequirePackage[letterpaper,margin=1in]{geometry}
\RequirePackage{amsmath,amsthm,mathtools}
\RequirePackage{fontspec,unicode-math}
\setmainfont{texgyretermes}[Extension=.otf,UprightFont=*-regular,BoldFont=*-bold,ItalicFont=*-italic,BoldItalicFont=*-bolditalic]
\setmathfont{texgyretermes-math.otf}
\RequirePackage{microtype,tikz-cd,enumitem,needspace,xcolor,booktabs,titlesec,etoolbox}
\definecolor{teal}{HTML}{176568}
\definecolor{burgundy}{HTML}{782F40}
\definecolor{inklink}{HTML}{176568}
\definecolor{accent}{HTML}{176568}
\titleformat{\section}{\large\bfseries\color{teal}}{\thesection}{.65em}{}
\titleformat{\subsection}{\normalsize\bfseries\color{burgundy}}{\thesubsection}{.65em}{}
\titleformat{\paragraph}[runin]{\normalsize\bfseries\color{burgundy}}{}{0pt}{}
\titlespacing*{\section}{0pt}{2.2ex plus .5ex minus .2ex}{1ex plus .2ex}
\titlespacing*{\subsection}{0pt}{1.6ex plus .4ex minus .2ex}{.7ex plus .2ex}
\RequirePackage{hyperref,bookmark}
\hypersetup{colorlinks=true,linkcolor=teal,urlcolor=teal,citecolor=teal,pdfstartview={XYZ null null 1.5},bookmarksnumbered=true}
\urlstyle{same}
\setcounter{tocdepth}{2}
\setlist[enumerate]{label=\arabic*.,ref=\arabic*,leftmargin=*,itemsep=4pt,topsep=4pt}
\setlength{\parindent}{1em}
\setlength{\parskip}{3pt}
\setlength{\emergencystretch}{2em}
\raggedbottom
\AddToHook{cmd/section/before}{\Needspace{8\baselineskip}}
\AddToHook{cmd/subsection/before}{\Needspace{6\baselineskip}}
\makeatletter
\newcommand{\afterblock}{\@afterindentfalse\@afterheading}
\makeatother
\newcounter{result}[section]
\renewcommand{\theresult}{\thesection.\arabic{result}}
\newcommand{\resulttitle}[3]{#1 #2\ifstrempty{#3}{}{ (#3)}.}
\newenvironment{result}[3]{\par\addvspace{7pt}\Needspace{4\baselineskip}\refstepcounter{result}\hypertarget{#2}{}\label{#2}\noindent{\color{burgundy}\hypersetup{linkcolor=burgundy}\hyperlink{proof:#2}{\textbf{\resulttitle{#1}{\theresult}{#3}}}}\quad\ignorespaces}{\par\addvspace{4pt}\afterblock}
\newenvironment{entry}[3]{\par\addvspace{7pt}\Needspace{3\baselineskip}\refstepcounter{result}\hypertarget{#2}{}\label{#2}\noindent{\color{burgundy}\textbf{\resulttitle{#1}{\theresult}{#3}}}\quad\ignorespaces}{\par\addvspace{4pt}\afterblock}
\newenvironment{demonstration}[1]{\par\noindent\hypertarget{proof:#1}{}{\color{burgundy}\hypersetup{linkcolor=burgundy}\hyperlink{#1}{\textbf{Proof.}}}\quad\ignorespaces}{\unskip\nobreak\hfill\qedsymbol\par\addvspace{5pt}\afterblock}
\newenvironment{citedresult}[4]{\par\addvspace{7pt}\Needspace{4\baselineskip}\refstepcounter{result}\hypertarget{#2}{}\label{#2}\noindent{\color{burgundy}\hypersetup{urlcolor=burgundy}\href{#4}{\textbf{\resulttitle{#1}{\theresult}{#3}}}}\quad\ignorespaces}{\par\addvspace{4pt}\afterblock}
\newcommand{\usedby}[1]{\par\noindent{\small\textbf{Used in:} #1.}\par\afterblock}
\newcommand{\coursebold}[1]{{\color{burgundy}\textbf{#1}}}
\newcommand{\coursetitle}[3]{\begin{center}{\large\bfseries\color{teal}#1}\\[4pt]{\LARGE\bfseries\color{teal}#2}\\[6pt]Math 615: Algebraic Topology I\\Fall 2026\quad\textbar\quad #3\end{center}}
\newcommand{\coursecontents}{{\small\setlength{\parskip}{0pt}\tableofcontents}}
\newcommand{\homeworkstyle}{\titleformat{\section}{\large\bfseries\color{teal}}{Exercise \thesection.}{.65em}{}}
% END SHARED COURSE STYLE
\newcommand{\Mor}{\operatorname{Mor}}\newcommand{\im}{\operatorname{im}}\newcommand{\coker}{\operatorname{coker}}\newcommand{\A}{\mathcal A}\newcommand{\Z}{\mathbb Z}\newcommand{\id}{\mathrm{id}}\DeclareMathOperator{\Cone}{Cone}\DeclareMathOperator{\Cyl}{Cyl}
\newcommand{\B}{\mathcal B}
\newcommand{\LF}{\mathrm L}
\newcommand{\RF}{\mathrm R}
\DeclareMathOperator{\Tor}{Tor}
\DeclareMathOperator{\Ext}{Ext}
\homeworkstyle
\hypersetup{pdftitle={Homework 3: Derived Functors, Tor, and Ext},pdfauthor={Math 615 Fall 2026},pdfsubject={Revision 3}}
\begin{document}
\coursetitle{Homework 3}{Derived Functors, Tor, and Ext}{September 14 and 16}
Let $\A$ and $\B$ be abelian categories. An object $P$ is \emph{projective} if
$\Mor_{\A}(P,-)$ takes epimorphisms to surjections; an object $I$ is
\emph{injective} if $\Mor_{\A}(-,I)$ takes monomorphisms to surjections.
The category $\A$ has \emph{enough projectives} if every object is a quotient
of a projective object, and \emph{enough injectives} if every object embeds
in an injective object.

Use homological grading for projective resolutions and cohomological grading
for injective resolutions. All functors between abelian categories below are
additive. You may use Homework 2, including homotopy invariance of homology and
the long exact sequence of a short exact sequence of complexes, as well as
their cochain analogues. In an abstract abelian category, use morphisms and
universal properties; element arguments are permitted for modules.

\section{Resolutions and Comparison Maps}\label{ex:comparison}
Assume that $\A$ has enough projectives. A \emph{projective resolution} of $M$
is an exact augmented complex
\[
\cdots\longrightarrow P_2\xrightarrow{d_2}P_1\xrightarrow{d_1}P_0
\xrightarrow{\epsilon}M\longrightarrow0
\]
with every $P_n$ projective.
\begin{enumerate}
\item Construct a projective resolution of every $M\in\A$ by successively
resolving kernels. For modules over a unital ring $R$, explain why free
modules suffice. Write an arbitrary free module as $\bigoplus_{i\in I}R$.
\item Given projective resolutions $P_\bullet\to M$ and $Q_\bullet\to N$,
prove that every $f:M\to N$ lifts to a chain map $\widetilde f:P_\bullet\to
Q_\bullet$. Prove that any two lifts of the same $f$ are chain homotopic.
Specify the induction that constructs both the lift and the homotopy.
\item Deduce that two projective resolutions of the same object are chain
homotopy equivalent by maps lifting its identity. Explain why the induced
maps on homology after applying any additive functor are independent of all
choices and compatible with composition.
\end{enumerate}

\section{Left Derived Functors}\label{ex:left}
Assume that $\A$ has enough projectives, and let $F:\A\to\B$ be right exact.
For a projective resolution $P_\bullet\to M$, set
\[
\LF_nF(M)=H_n(F(P_\bullet)),\qquad n\geq0.
\]
\begin{enumerate}
\item Using Exercise~\ref{ex:comparison}, define $\LF_nF$ on morphisms and
prove that it is an additive functor, well defined up to canonical natural
isomorphism independently of the chosen resolutions.
\item Construct a natural isomorphism $\LF_0F\cong F$. Identify precisely
where right exactness is used.
\item Prove that $\LF_nF(P)=0$ for every projective $P$ and $n>0$.
If $F$ is exact, prove that $\LF_nF(M)=0$ for every $M$ and $n>0$.
\end{enumerate}


\section{The Horseshoe Lemma and Connecting Maps}\label{ex:horseshoe}
Assume that $\A$ has enough projectives. Let
$0\to M'\to M\to M''\to0$ be exact, and choose projective resolutions
$P'_\bullet\to M'$ and $P''_\bullet\to M''$.
\begin{enumerate}
\item Construct a projective resolution $P_\bullet\to M$ and a short exact
sequence of augmented complexes lifting the given sequence such that
\[
0\longrightarrow P'_n\longrightarrow P_n\longrightarrow P''_n
\longrightarrow0
\]
is split exact for every $n\geq0$. You may take $P_n=P'_n\oplus P''_n$;
construct its differential and augmentation. The given sequence of objects
is not assumed to split.
\item For a right exact $F:\A\to\B$, explain why applying $F$ gives a short
exact sequence of the unaugmented complexes, even though $F$ need not be
left exact. Deduce a long exact sequence ending in
\[
\cdots\to\LF_1F(M)\to\LF_1F(M'')\xrightarrow{\partial}
F(M')\to F(M)\to F(M'')\to0.
\]
Describe the connecting maps in all positive degrees.
\item Prove that this long exact sequence is independent of the horseshoe
choices and natural in morphisms of short exact sequences. Explain how to
choose comparison maps compatible with the inclusions and projections.
\end{enumerate}

\section{Right Derived Functors}\label{ex:right}
Now assume that $\A$ has enough injectives and let $G:\A\to\B$ be left exact.
An \emph{injective resolution} of $M$ is an exact sequence
$0\to M\to I^0\to I^1\to\cdots$ with each $I^n$ injective. Set
\[
\RF^nG(M)=H^n(G(I^\bullet)),\qquad n\geq0.
\]
\begin{enumerate}
\item State and prove the injective analogues of the existence and comparison
results in Exercise~\ref{ex:comparison}. Give the direction of the
comparison map associated to $f:M\to N$. Deduce that $\RF^nG$ is a
well-defined additive functor up to canonical natural isomorphism.
\item Prove that $\RF^0G\cong G$, and that $\RF^nG(I)=0$ for injective $I$
and $n>0$. Construct the long exact sequence beginning in
\[
0\to G(M')\to G(M)\to G(M'')\xrightarrow{\delta} \RF^1G(M')
\to\RF^1G(M)\to\RF^1G(M'')\to\cdots.
\]
You may obtain the injective horseshoe lemma by duality.
\item Show that $G$ is exact if and only if $\RF^1G(M)=0$ for every $M$.
Prove the corresponding criterion for a right exact $F$ in terms of
$\LF_1F$, assuming enough projectives.
\end{enumerate}


\section{Dimension Shifting and Acyclic Resolutions}\label{ex:shift}
Continue with $F$ as in Exercise~\ref{ex:left} and $G$ as in
Exercise~\ref{ex:right}, making the corresponding assumptions on $\A$.
Call $A$ \emph{$F$-acyclic} if $\LF_nF(A)=0$ for every $n>0$.
\begin{enumerate}
\item If $0\to K\to P\to M\to0$ is exact and $P$ is projective, construct
natural isomorphisms
\[
\LF_{n+1}F(M)\cong\LF_nF(K)\quad(n\geq1)
\]
and the exact sequence
$0\to\LF_1F(M)\to F(K)\to F(P)\to F(M)\to0$.
Explain why the displayed isomorphism is not asserted for $n=0$.
\item Give and prove the dual statements for an exact sequence
$0\to M\to I\to C\to0$ with $I$ injective.
\item Suppose $\cdots\to A_1\to A_0\to M\to0$ is an exact resolution
by $F$-acyclic objects. Prove that
$H_n(F(A_\bullet))\cong\LF_nF(M)$ for every $n\geq0$.
Use the short exact sequences involving its successive kernels and
dimension shifting. No finite-length assumption is made.
\end{enumerate}

\section{Tensor Products and Tor}\label{ex:tor}
Let $R$ be a commutative unital ring and $N$ an $R$-module. Define
\[
\Tor_n^R(M,N)=\LF_n(-\otimes_R N)(M).
\]
An $R$-module $N$ is \emph{flat} if $-\otimes_R N$ is exact.
\begin{enumerate}
\item Check that $-\otimes_R N$ is right exact. Use
Exercise~\ref{ex:right} to prove that $N$ is flat if and only if
$\Tor_1^R(M,N)=0$ for every $M$.
\item Let $r\in R$ be a non-zero-divisor. Use the resolution
$0\to R\xrightarrow{r}R\to R/(r)\to0$ to compute
$\Tor_n^R(R/(r),N)$ for all $n\geq0$. Express the answer using $N/rN$
and $N[r]=\{x\in N:rx=0\}$. Explain why the same resolution cannot
be used when $r$ is a zero-divisor.
\item Take $R=\Z$ and positive integers $m,n\geq2$. Compute
$\Tor_i^{\Z}(\Z/m,\Z/n)$ for every $i\geq0$.
Identify the connecting injection
$\Tor_1^{\Z}(\Z/m,N)\to N$ obtained by tensoring
$0\to\Z\xrightarrow{m}\Z\to\Z/m\to0$.
\end{enumerate}


\section{Ext and Extensions}\label{ex:ext}
Let $R$ be a unital ring, not necessarily commutative, and work with left
$R$-modules. Fix a module $M$ and define
\[
\Ext_R^n(M,N)=\RF^n\Mor_R(M,-)(N).
\]
You may use that module categories have enough injectives and the standard
comparison theorem identifying these groups with
$H^n(\Mor_R(P_\bullet,N))$ for a projective resolution $P_\bullet\to M$.
Here the cochain differential sends $u:P_n\to N$ to $u\circ d_{n+1}$.
For a noncommutative $R$, these are abelian groups, not in general left
$R$-modules.
\begin{enumerate}
\item Given $0\to K\xrightarrow{j}P\to M\to0$ with $P$ projective,
prove that
\[
\Ext_R^1(M,N)\cong
\coker\bigl(\Mor_R(P,N)\xrightarrow{j^*}\Mor_R(K,N)\bigr).
\]
For $u:K\to N$, form the pushout along $u$ and obtain an extension
$0\to N\to E_u\to M\to0$. Prove that the resulting construction
identifies the displayed cokernel with equivalence classes of extensions,
where equivalences induce the identities on $N$ and $M$.
Show that the zero class corresponds exactly to split extensions.
\item Deduce that $M$ is projective if and only if
$\Ext_R^1(M,N)=0$ for every left $R$-module $N$.
\item For $R=\Z$, $m\geq2$, and an abelian group $N$, compute
$\Ext_{\Z}^i(\Z/m,N)$ in every degree. Then take $N=\Z/n$ and express
the answer in terms of $\gcd(m,n)$.
\end{enumerate}

\section{An Infinite Resolution}\label{ex:periodic}
Let $k$ be a field, let $R=k[\varepsilon]/(\varepsilon^2)$, and regard
$k=R/(\varepsilon)$ as an $R$-module.
\begin{enumerate}
\item Prove that the augmented complex
\[
\cdots\xrightarrow{\varepsilon}R\xrightarrow{\varepsilon}R
\xrightarrow{\varepsilon}R\longrightarrow k\longrightarrow0
\]
is a free resolution of $k$ by computing its kernels and images.
\item Compute $\Tor_n^R(k,k)$ and $\Ext_R^n(k,k)$ for every $n\geq0$,
including the differentials in the complexes used for these computations.
\item Prove that $k$ has no projective resolution of finite length over $R$.
Contrast this with the computations over $\Z$ in
Exercises~\ref{ex:tor} and~\ref{ex:ext}.
\end{enumerate}
\end{document}
